% Task 4 report, built from reports/template/report_template.tex
% build from this folder with: tectonic task4_report.tex
\documentclass[11pt,a4paper]{article}
\usepackage[margin=2cm,top=1.4cm]{geometry}
\usepackage{enumitem}
\usepackage{hyperref}

% title block: left-aligned title, then author on the left and date on the right in one line
\makeatletter
\renewcommand{\maketitle}{%
  \noindent{\LARGE \@title\par}\vspace{0.6em}
  \noindent{\large \@author\hfill\@date\par}\vspace{3em}}
\makeatother

\title{Task 4: Appearance-Aware Denoising Module}
\author{L335 MiniProject}
\date{\today}

\begin{document}
\maketitle

% task description, copied verbatim from the project brief and split into bullet points
% dashes at 1em, item text at 2.5em; no extra space before, between or after the items
\emergencystretch=2em% lets TeX stretch spaces a little more instead of pushing long compounds into the margin
\noindent\textbf{Task description:}
\begin{itemize}[label=--, leftmargin=2.5em, rightmargin=2.5em, labelindent=1em, labelwidth=1em, labelsep=0.5em, align=left, topsep=0pt, partopsep=0pt, itemsep=0pt, parsep=0pt]
\item Propose and implement an appearance-aware denoising module.
\item Extend one pretrained baseline with a small colour-conditioned residual module or neighbourhood-weighting mechanism.
\item Compare geometry-only features, straightforward RGB fusion, and a learned gate controlling how strongly colour affects the correction.
\item Keep the training data, module capacity, and training budget comparable.
\end{itemize}
% optional note below the list, aligned with the dashes:
% \noindent\hspace*{1em}\textit{Note:} ...
\vspace{1.3em}% gap before the first section; \section adds about 1.5em itself and the larger of the two wins

\end{document}
